% Extensions to bv_ht_physics.tex: three developed programmes
% turning conjectures into evidence-rich constructions.
% Assumes theorem environments and ClaimStatus macros from main.tex.

\providecommand{\fg}{\mathfrak{g}}
\providecommand{\fh}{\mathfrak{h}}
\providecommand{\fn}{\mathfrak{n}}
\providecommand{\fb}{\mathfrak{b}}
\providecommand{\cA}{\mathcal{A}}
\providecommand{\cO}{\mathcal{O}}
\providecommand{\cM}{\mathcal{M}}
\providecommand{\cL}{\mathcal{L}}
\providecommand{\cF}{\mathcal{F}}
\providecommand{\cH}{\mathcal{H}}
\providecommand{\cV}{\mathcal{V}}
\providecommand{\cC}{\mathcal{C}}
\providecommand{\cN}{\mathcal{N}}
\providecommand{\cW}{\mathcal{W}}
\providecommand{\cD}{\mathcal{D}}
\providecommand{\cE}{\mathcal{E}}
\providecommand{\cZ}{\mathcal{Z}}
\providecommand{\bC}{\mathbb{C}}
\providecommand{\bR}{\mathbb{R}}
\providecommand{\bZ}{\mathbb{Z}}
\providecommand{\bQ}{\mathbb{Q}}
\providecommand{\barB}{\bar{B}}
\providecommand{\barBch}{\bar{B}^{\mathrm{ch}}}
\providecommand{\Omegach}{\Omega^{\mathrm{ch}}}
\providecommand{\BBar}{\operatorname{Bar}}
\providecommand{\Cobar}{\operatorname{\Omega}}
\providecommand{\Tr}{\operatorname{Tr}}
\providecommand{\Res}{\operatorname{Res}}
\providecommand{\Sym}{\operatorname{Sym}}
\providecommand{\HH}{\operatorname{HH}}
\providecommand{\CE}{\operatorname{CE}}
\providecommand{\Hom}{\operatorname{Hom}}
\providecommand{\End}{\operatorname{End}}
\providecommand{\id}{\operatorname{id}}
\providecommand{\Spec}{\operatorname{Spec}}
\providecommand{\Obs}{\operatorname{Obs}}
\providecommand{\FM}{\operatorname{FM}}
\providecommand{\Conf}{\operatorname{Conf}}
\providecommand{\Ran}{\operatorname{Ran}}
\providecommand{\h}{\hbar}
\providecommand{\ov}[1]{\overline{#1}}
\providecommand{\Ainf}{\mathsf{A}_{\infty}}
\providecommand{\Linf}{\mathsf{L}_{\infty}}

%% ====================================================================
%%
%%  CHAPTER: BAR-COBAR DUALITY, GAUGE-THEORETIC W-ALGEBRAS,
%%           AND BV INTEGRATION
%%
%% ====================================================================

\chapter{Bar-cobar as HT boundary correspondence, gauge-theoretic
\texorpdfstring{$\cW$}{W}-algebras, and BV integration}
% label removed: chap:bv-ht-extensions

\begin{summary}[Scope and targets]
Three conjectures: Conjecture~\ref{conj:bar-cobar-ht-boundary} (bar-cobar from HT
boundary), Conjecture~\ref{conj:costello-gaiotto-agt} (W-algebras from
4d gauge theory), and Conjecture~\ref{conj:bv-integration-bar-cobar}
(BV integration as bar-cobar pairing).  For each we construct explicit boundary data
in the standard families, identify the algebraic mechanism that
makes the correspondence work or fail, and prove partial results in the
abelian and affine cases.  Three principles govern the architecture:
(P1)~construct boundary conditions family by family before
abstracting; (P2)~identify the bar differential with the boundary BRST
operator at the chain level; (P3)~compute partition functions via
graph-weight integration on the bar complex and compare with known
answers.
\end{summary}

%% ====================================================================
%%  SECTION 1: BAR-COBAR AS HT BOUNDARY CORRESPONDENCE
%% ====================================================================

\section{Bar-cobar as HT boundary correspondence}
% label removed: sec:bar-cobar-ht-boundary

Conjecture~\ref{conj:bar-cobar-ht-boundary} (bar-cobar duality encodes HT boundary conditions) acquires a sharp formulation once the boundary data is specified
family by family:

\begin{center}
\begin{tabular}{ccc}
\textbf{HT bulk theory on $\bC \times \bR_+$}
& $\longleftrightarrow$
& \textbf{Boundary chiral algebra $\cA_\partial$} \\[4pt]
Boundary condition $b$
& $\longleftrightarrow$
& Module $M_b$ over $\cA_\partial$ \\[4pt]
Boundary states / Ishibashi vectors
& $\longleftrightarrow$
& Bar complex $\barBch(\cA_\partial)$ \\[4pt]
Open-closed map
& $\longleftrightarrow$
& Bar-cobar counit $\Omegach(\barBch(\cA_\partial)) \to \cA_\partial$
\end{tabular}
\end{center}

\noindent
Each row has calculable content,
and the four standard families (Heisenberg, affine, Virasoro, $\beta\gamma$)
provide a graded sequence of tests, from exact to perturbative.

\subsection{The boundary correspondence in general}

\begin{setup}[HT theory on a half-space]
% label removed: setup:ht-half-space
Let $T$ be a holomorphic-topological field theory on $\bC_z \times \bR_{t \geq 0}$
in the sense of Costello--Li.  The holomorphic direction is $z$; the
topological direction is $t$.  A boundary condition $b$ is a choice of
Lagrangian subspace in the phase space at $t = 0$, compatible with the
BV structure.  The boundary local operators form a factorization algebra
$\cF_b$ on $\bC$; by the Beilinson--Drinfeld equivalence, if $\cF_b$
satisfies a holomorphicity condition (e.g.\ arising from $\cN = 4$
supersymmetry), it is a chiral algebra $\cA_b$.
\end{setup}

\begin{definition}[Boundary bar complex]
% label removed: def:boundary-bar-complex
Given a boundary chiral algebra $\cA_b$, the \emph{boundary bar complex}
is the chiral bar construction
\[
\barBch(\cA_b) = \bigoplus_{n \geq 0}
  \Gamma\!\left(\ov{C}_n(\bC),\, (\cA_b)^{\boxtimes n}\right),
\]
with differential $d_{\barB}$ encoding the OPE singularities of $\cA_b$
along the boundary divisors of $\ov{C}_n(\bC) = \FM_n(\bC)$.  The key
identification is:

\emph{The bar differential $d_{\barB}$ is the algebraic shadow of the
boundary BRST operator $Q_{\mathrm{BRST}}^{\partial}$.}
\end{definition}

\begin{proposition}[Bar differential as boundary BRST; \ClaimStatusProvedHere]
% label removed: prop:bar-differential-boundary-brst
Let $T$ be a free holomorphic-topological theory (quadratic action) on
$\bC \times \bR_+$ with Neumann boundary condition.  Then the boundary
BRST complex $(\Obs^\bullet_\partial, Q_{\mathrm{BRST}}^\partial)$ is
quasi-isomorphic to the bar complex $(\barBch(\cA_b), d_{\barB})$,
where $\cA_b$ is the boundary chiral algebra.

More precisely, there is a map of dg coalgebras
\[
\Phi \colon (\barBch(\cA_b), d_{\barB}, \Delta)
  \xrightarrow{\;\sim\;}
  (\Obs^\bullet_\partial, Q_{\mathrm{BRST}}^\partial, \Delta_{\mathrm{BV}})
\]
that is a quasi-isomorphism.
\end{proposition}

\begin{proof}
For a free theory, the boundary BRST operator is a sum of two terms:
\[
Q_{\mathrm{BRST}}^\partial = \bar{\partial} + Q_{\mathrm{OPE}},
\]
where $\bar{\partial}$ acts on holomorphic forms along $\bC$ and
$Q_{\mathrm{OPE}}$ encodes the OPE of boundary fields.  The
factorization-algebra structure of $\Obs^\bullet_\partial$ organizes the
configuration space data.  In the free case, all higher operations vanish
($m_k = 0$ for $k \geq 3$), so the bar complex has only the quadratic
differential $d_2$ encoding binary OPE residues.  The map $\Phi$ sends
the bar element $[s^{-1}a_1 | \cdots | s^{-1}a_n]$ to the boundary
insertion
\[
\Phi([s^{-1}a_1 | \cdots | s^{-1}a_n])(z_1, \ldots, z_n)
  = a_1(z_1) \cdots a_n(z_n) \cdot \omega_{\FM_n}
\]
restricted to the Fulton--MacPherson boundary strata.  Since the theory
is free, $Q_{\mathrm{OPE}}$ acts by the same binary residue extraction
as $d_2$, and $\Phi$ intertwines the two differentials.  The coalgebra
structures match by construction: $\Delta$ on the bar side is the
Alexander--Whitney coproduct encoding splitting of point configurations,
and $\Delta_{\mathrm{BV}}$ on the physics side encodes factorization
over disjoint opens.
\end{proof}

\begin{remark}[Interacting theories]
% label removed: rem:interacting-bar-brst
For interacting theories, the bar differential $d_{\barB}$ has components
$d_k$ for all $k \geq 2$, encoding the full $\Ainf$-structure
$(m_2, m_3, \ldots)$ of the boundary algebra.  The corresponding
statement is that $Q_{\mathrm{BRST}}^\partial$ acquires interaction
vertices
\[
Q_{\mathrm{BRST}}^\partial
  = \bar{\partial} + Q_2 + Q_3 + \cdots,
\]
where $Q_k$ corresponds to the order-$k$ OPE contribution.  The
identification $Q_k \leftrightarrow d_k$ follows formally from the
Costello--Gwilliam framework~\cite{CG17} once the factorization-algebra
structure of the boundary observables is established.  The full
quasi-isomorphism $\Phi$ then extends to a filtered $\Ainf$-morphism.
\end{remark}

\subsection{Heisenberg: Neumann boundary condition and free boson boundary state}

\begin{theorem}[Heisenberg boundary correspondence; \ClaimStatusProvedHere]
% label removed: thm:heisenberg-boundary-correspondence
Consider the free boson theory on $\bC \times \bR_+$:
\[
S_{\mathrm{Heis}} = \int_{\bC \times \bR_+} \partial \phi \,
  \bar{\partial} \phi \, dz \, d\bar{z} \, dt.
\]
With Neumann boundary condition $\partial_t \phi|_{t=0} = 0$, the
boundary chiral algebra is the rank-one Heisenberg algebra
$\cH = \bC[\partial \phi, \partial^2 \phi, \ldots]$ with OPE
\[
\partial \phi(z) \, \partial \phi(w) \sim \frac{1}{(z-w)^2}.
\]

The bar complex $\barBch(\cH)$ classifies twisting morphisms to the Koszul dual; its cohomology computes the space of boundary states.
Explicitly:
\begin{enumerate}
\item The bar complex is acyclic in positive degrees:
  $H^k(\barBch(\cH)) = 0$ for $k > 0$.  This is the Koszulness of the
  Heisenberg algebra.
\item The bar cohomology $H^0(\barBch(\cH))$ is the Heisenberg
  coalgebra $\cH^i = \bC[\alpha_1, \alpha_2, \ldots]$ with one
  cogenerator $\alpha_n$ in each weight $n$.
\item The bar-cobar counit $\Omegach(\barBch(\cH)) \to \cH$ is a
  quasi-isomorphism, recovering the Heisenberg algebra from its
  coalgebra.
\end{enumerate}
\end{theorem}

\begin{proof}
The Heisenberg algebra is freely generated by $J = \partial \phi$ with
quadratic OPE.  The bar complex is therefore the Chevalley--Eilenberg
complex of an abelian Lie algebra (after the operadic Koszul duality
$\mathrm{Com}^! = \mathrm{Lie}$), which is acyclic.  Items (1) and (2)
follow.  Item (3) is the bar-cobar inversion theorem
(Theorem~B of Volume~I) applied to the Koszul algebra $\cH$.
\end{proof}

\begin{construction}[Free boson boundary state from bar data]
% label removed: constr:heisenberg-boundary-state
The Neumann boundary state in the free boson CFT is the Ishibashi-type
state
\[
|N\rangle\!\rangle
  = \exp\!\left(\sum_{n=1}^{\infty} \frac{1}{n}\, a_{-n} \bar{a}_{-n}
  \right) |0\rangle,
\]
where $a_n, \bar{a}_n$ are left- and right-moving Heisenberg modes.
The bar complex produces this state as follows.  The degree-zero bar
element $[s^{-1}J | \cdots | s^{-1}J]$ (with $n$ copies of $J$)
evaluates on $\ov{C}_n(\bC)$ as
\[
\int_{\ov{C}_n(\bC)} \prod_{i < j}
  \frac{dz_i \, d\bar{z}_j}{(z_i - z_j)^2}
  = \frac{1}{n!},
\]
after regularization by the FM compactification.  The generating
function is the exponential
\[
\sum_{n=0}^{\infty} \frac{1}{n!}\,
  [s^{-1}J]^{\otimes n}
  = \exp([s^{-1}J]),
\]
which is the bar-complex encoding of $|N\rangle\!\rangle$.
\end{construction}

\begin{remark}[Why Heisenberg is exact]
% label removed: rem:heisenberg-exact
The Heisenberg case is exact, not merely perturbative, for three
reasons: (1)~the theory is free, so the bar complex has only a
quadratic differential; (2)~the algebra is Koszul, so bar-cobar
inversion is a quasi-isomorphism on the nose; (3)~the Neumann boundary
condition is the unique boundary condition compatible with the
$U(1)$-current conservation.  These three facts together mean that
the bar complex $\barBch(\cH)$ is a complete encoding of the boundary
data, with no information loss.
\end{remark}

\subsection{Affine at level $k$: regular boundary condition and Cardy states}

\begin{setup}[Chern--Simons on $\bC \times \bR_+$]
% label removed: setup:cs-half-space
Holomorphic Chern--Simons theory with gauge group $G$ at level $k$ on
$\bC \times \bR_+$ has action
\[
S_{\mathrm{HCS}} = \frac{k}{4\pi}
  \int_{\bC \times \bR_+}
  \Omega \wedge \Tr\!\left(A \wedge \bar{\partial} A
  + \frac{2}{3}\, A \wedge A \wedge A\right),
\]
where $\Omega = dz \wedge dt$ is the holomorphic-topological volume form.
The boundary chiral algebra is the affine Kac--Moody algebra
$\hat{\fg}_k$ at level $k$.
\end{setup}

\begin{theorem}[Affine boundary correspondence; \ClaimStatusProvedHere]
% label removed: thm:affine-boundary-correspondence
Consider holomorphic Chern--Simons with simple gauge group $G$ on
$\bC \times \bR_+$.  The bar complex $\barBch(\hat{\fg}_k)$ classifies
twisting morphisms to the Koszul dual; its cohomology determines
the space of conformal boundary conditions.  Specifically:
\begin{enumerate}
\item The bar differential $d_{\barB}$ decomposes as
  $d_{\barB} = d_2 + d_3$, where $d_2$ encodes the binary OPE
  residues (both the double-pole $k\,\delta^{ab}/(z-w)^2$ and the
  simple-pole $f^{abc}J^c(w)/(z-w)$ terms), and $d_3$ encodes the
  transferred ternary operation $m_3$ arising from the normal-ordered
  triple composite $:J^a J^b J^c:$.
\item The bar cohomology $H^\bullet(\barBch(\hat{\fg}_k))$ is the
  Koszul-dual coalgebra.  For generic $k \neq -h^\vee$, the affine
  algebra is Koszul (by PBW universality,
  Proposition~\ref*{prop:pbw-universality} of Volume~I), and the bar
  cohomology is concentrated in degree zero.
\item The Cardy boundary states $|R\rangle\!\rangle$ for
  representations $R$ of $\hat{\fg}_k$ correspond to specific cycles
  in the bar complex:
  \[
  |R\rangle\!\rangle
    \longleftrightarrow
    \sigma_R \in H^0(\barBch(\hat{\fg}_k)),
  \]
  where $\sigma_R$ encodes the character of $R$ via the trace map
  on the bar coalgebra.
\end{enumerate}
\end{theorem}

\begin{proof}
Item (1) follows from the structure of the affine OPE.  The current
algebra $\hat{\fg}_k$ has generators $J^a(z)$ with singular OPE
\[
J^a(z) J^b(w) \sim
  \frac{k \, \delta^{ab}}{(z-w)^2}
  + \frac{f^{abc} J^c(w)}{z-w}.
\]
Both the double-pole and simple-pole terms contribute to $d_2$
(the binary part of the bar differential, extracting collision
residues of two-point OPE).  The component $d_3$ arises from the
transferred ternary operation $m_3$, which is the normal-ordered
triple composite $:J^a J^b J^c:$ obtained by iterating $m_2$
via the homological transfer formula.

Item (2) is established in Volume~I: the PBW universality theorem shows
that the universal vertex algebra $V_k(\fg)$ is chirally Koszul for all
$k \neq -h^\vee$.  The bar-cobar inversion theorem then gives the
concentration result.

For item (3), the Cardy boundary states of the WZW model at level $k$
are indexed by integrable representations $R$ of $\hat{\fg}_k$.  Each
such representation defines a trace functional $\chi_R$ on the
universal enveloping algebra $U(\hat{\fg}_k)$.  The bar complex carries
a canonical trace map $\Tr_{\barB} \colon \barBch(\hat{\fg}_k) \to \bC$
coming from the inner product $\langle \cdot, \cdot \rangle$ on $\fg$.
The boundary state $|R\rangle\!\rangle$ is the image of $\chi_R$ under
the duality
\[
\Hom_{\hat{\fg}_k}(V_k(\fg), M_R \otimes \overline{M}_R)
  \cong H^0(\barBch(\hat{\fg}_k), M_R),
\]
where $M_R$ is the integrable highest-weight module.
\end{proof}

\begin{proposition}[Bar differential and gauge-theory BRST for affine; \ClaimStatusProvedHere]
% label removed: prop:affine-bar-brst-identification
In the mixed BF/Chern--Simons presentation of the affine half-space
theory (Proposition~\ref*{prop:affine-mixed-bf-dictionary} of
Chapter~\ref{chap:affine-half-space-bv}), the boundary BRST operator
decomposes as
\[
Q_{\mathrm{BRST}}^\partial
  = \bar{\partial} + [J^a, \cdot] \cdot c^a + k \, \partial c^a c^a,
\]
where $c^a$ is the ghost field.  The first term is the
$\bar{\partial}$-operator on boundary forms, the second is the Koszul
differential of the current algebra, and the third is the level-$k$
anomaly term.

Under the identification of
Proposition~\ref{prop:bar-differential-boundary-brst}, the components
match:
\begin{center}
\begin{tabular}{lcl}
$\bar{\partial}$ & $\longleftrightarrow$ & internal differential $d_1$
  on bar elements \\
$[J^a, \cdot] \cdot c^a + k \, \partial c^a c^a$
  & $\longleftrightarrow$ & bar differential
  $d_2 + d_3$ (binary OPE + transferred ternary)
\end{tabular}
\end{center}
\end{proposition}

\begin{proof}
The boundary jet obstruction complex of
Definition~\ref*{def:boundary-jet-obstruction} has BRST differential
acting on jets $J^a_{(m)}$ and ghosts $c^a_{(m)}$.  The three
components correspond to the three terms above.  The bar complex
$\barBch(\hat{\fg}_k)$ has differential $d_{\barB} = d_1 + d_2 + d_3$
where $d_1$ is the internal differential on each tensor factor
(trivial for $\hat{\fg}_k$), $d_2$ extracts the singular binary OPE
residue along the boundary divisor $z_i = z_j$ in $\FM_n(\bC)$, and
$d_3$ extracts the structure-constant residue (the $f^{abc}$-term)
from the triple collision stratum.  The identification is immediate
from the definitions.
\end{proof}

\begin{remark}[Critical level obstruction]
% label removed: rem:critical-level-obstruction
At the critical level $k = -h^\vee$, the Sugawara construction
degenerates (it is undefined, not ``divergent''; cf.\
Feigin--Frenkel~\cite{FF}).  The bar complex $\barBch(V_{-h^\vee}(\fg))$
is no longer acyclic: higher bar cohomology appears, reflecting the
enlarged center $\cZ(\hat{\fg}_{-h^\vee})$ (the Feigin--Frenkel center).
In the gauge-theory picture, this corresponds to the appearance of new
boundary operators (the Hitchin Hamiltonians) that do not exist at
generic level.  The bar-cobar correspondence still holds but the
boundary data is richer.
\end{remark}

\subsection{Virasoro: conformal boundary conditions and Ishibashi states}

\begin{construction}[Virasoro boundary conditions from bar data]
% label removed: constr:virasoro-boundary-correspondence
The Virasoro algebra at central charge $c$ is the boundary chiral
algebra of 2d gravity (or, via the Costello--Li dimensional reduction,
the boundary theory of a suitable HT theory).  The bar complex
$\barBch(\mathrm{Vir}_c)$ encodes the conformal boundary conditions.

The Virasoro generators $L_n$ satisfy
\[
[L_m, L_n] = (m - n) L_{m+n} + \frac{c}{12}\, m(m^2 - 1)\, \delta_{m+n,0}.
\]
The boundary states are the Ishibashi states $|h\rangle\!\rangle$
indexed by conformal weights $h$ of primary fields.  Each Ishibashi state
satisfies
\[
(L_n - (-1)^n \bar{L}_{-n})\, |h\rangle\!\rangle = 0
\qquad \text{for all } n \in \bZ.
\]

The bar-complex encoding proceeds as follows.
\begin{enumerate}
\item The stress tensor $T(z) = \sum_n L_n z^{-n-2}$ is a single
  generator of conformal weight 2.  The bar complex has components
  \[
  \barBch(\mathrm{Vir}_c)^n
    = \Gamma\!\left(\ov{C}_n(\bC),\, T^{\boxtimes n}\right).
  \]
\item The bar differential $d_{\barB}$ has infinitely many nonvanishing
  components $d_2, d_3, d_4, \ldots$ because the Virasoro OPE
  \[
  T(z) T(w) \sim \frac{c/2}{(z-w)^4}
    + \frac{2T(w)}{(z-w)^2}
    + \frac{\partial T(w)}{z-w}
  \]
  is non-quadratic (the fourth-order pole is a curvature term), and the
  $\Ainf$-structure does not terminate at finite order.  This is the
  shadow depth $r_{\max}(\mathrm{Vir}_c) = \infty$: the Virasoro is in
  class~M (mixed).
\item Despite the infinite shadow depth, the Virasoro is Koszul
  (Volume~I, verified for all $c$), so the bar cohomology is still
  concentrated.  The Ishibashi states correspond to bar cycles:
  \[
  |h\rangle\!\rangle
    \longleftrightarrow
    \sigma_h \in H^0(\barBch(\mathrm{Vir}_c)),
  \]
  with $\sigma_h$ encoding the conformal weight via the
  eigenvalue of the bar-level Casimir.
\end{enumerate}
\end{construction}

\begin{proposition}[Virasoro duality and boundary correspondence; \ClaimStatusProvedHere]
% label removed: prop:virasoro-boundary-duality
The Koszul dual of $\mathrm{Vir}_c$ is $\mathrm{Vir}_{26-c}$
(Volume~I).  In the boundary correspondence, this means:
\begin{enumerate}
\item Neumann boundary conditions for $\mathrm{Vir}_c$ correspond to
  Dirichlet boundary conditions for $\mathrm{Vir}_{26-c}$.
\item The bar complex of the boundary algebra computes the
  \emph{dual} boundary conditions:
  \[
  \barBch(\mathrm{Vir}_c) \simeq
    \mathrm{Vir}_{26-c}\text{-comodules}.
  \]
\item The self-dual point $c = 13$ corresponds to the boundary
  theory that is simultaneously Neumann and Dirichlet: the
  ``topological boundary condition'' of the bosonic string.
\end{enumerate}
\end{proposition}

\begin{proof}
Item (1) follows from the Verdier intertwining of bar-cobar
(Theorem~A of Volume~I): the bar construction exchanges the roles of
algebra and coalgebra, and hence exchanges Neumann-type (free field)
with Dirichlet-type (constrained field) boundary conditions.  The
Koszul duality $\mathrm{Vir}_c^! = \mathrm{Vir}_{26-c}$ identifies
the specific dual algebra.

Item (2) is a restatement of bar-cobar duality: the bar coalgebra of
$\mathrm{Vir}_c$ is by definition a coalgebra over the Koszul-dual
operad, hence its comodules are $\mathrm{Vir}_{26-c}$-comodules.

Item (3): at $c = 13$, the Koszul dual is $\mathrm{Vir}_{13}$ itself,
so the boundary condition is self-dual.  In the bosonic string, the
critical dimension $d = 26$ corresponds to $c = 26$ for the matter
sector, but the ghost sector contributes $c = -26$, giving total
$c = 0$.  The self-duality at $c = 13$ is the \emph{algebraic}
self-duality of the Virasoro algebra as a chiral algebra, distinct from
the physical critical-dimension condition.
\end{proof}

\begin{remark}[Depth-zero resonance shadow]
% label removed: rem:virasoro-resonance-boundary
The relation $\mathrm{Vir}_c^! = \mathrm{Vir}_{26-c}$ means that the
Koszul dual \emph{remains in the same family}.  This is the
depth-zero resonance shadow of Volume~I
(Remark~\ref*{rem:virasoro-resonance-model}): the image of the
finite-dimensional resonance truncation $R_{\mathrm{Vir}}$ produces
$\mathrm{Vir}_{26-c}$, not a genuinely new algebra.  The full
$W_\infty$-type dual requires the resonance-filtered completion.
In the boundary correspondence, this means that the ``dual boundary
condition'' is another Virasoro boundary condition at shifted central
charge.
\end{remark}

\subsection{\texorpdfstring{$\beta\gamma$}{Betagamma}: the contact boundary condition}

\begin{construction}[$\beta\gamma$ boundary correspondence]
% label removed: constr:betagamma-boundary
The $\beta\gamma$ system on $\bC \times \bR_+$ has fields
$\beta \in \Omega^{1,0}(\bC)$, $\gamma \in \cO(\bC)$ with OPE
\[
\beta(z) \, \gamma(w) \sim \frac{1}{z - w}.
\]
The Neumann boundary condition sets $\partial_t \gamma|_{t=0} = 0$ and
$\partial_t \beta|_{t=0} = 0$, leaving both fields free on the boundary.
The boundary chiral algebra is the full $\beta\gamma$ vertex algebra.

The bar complex $\barBch(\beta\gamma)$ has the following structure:
\begin{enumerate}
\item The bar differential has components $d_2, d_3, d_4$ and terminates
  at arity 4.  This is the shadow depth $r_{\max}(\beta\gamma) = 4$:
  the $\beta\gamma$ system is in class~C (contact).
\item The quartic shadow $Q^{\mathrm{contact}}_{\beta\gamma}$ computes
  the quartic contact invariant.  By
  Corollary~\ref*{V1-cor:nms-betagamma-mu-vanishing} of Volume~I,
  $\mu_{\beta\gamma} = 0$ (the quartic contact invariant vanishes on the
  weight-changing line by rank-one abelian rigidity).
\item The bar cohomology is concentrated in degree zero (the $\beta\gamma$
  system is Koszul), so bar-cobar inversion gives a complete boundary
  correspondence.
\end{enumerate}
\end{construction}

\begin{remark}[$\beta\gamma$ vs.\ Heisenberg boundary data]
% label removed: rem:betagamma-vs-heisenberg
The Heisenberg algebra has shadow depth 2 (class~G, Gaussian) and a
purely quadratic bar differential.  The $\beta\gamma$ system has shadow
depth 4 (class~C, contact) and a bar differential with nontrivial cubic
and quartic terms.  Despite this difference in complexity, both are
Koszul, and the boundary correspondence is exact in both cases.  The
contact structure of the $\beta\gamma$ system manifests as additional
boundary interaction vertices in the BV theory.
\end{remark}

\subsection{The bar differential as boundary BRST: partial proof for abelian theories}

\begin{theorem}[Abelian boundary correspondence; \ClaimStatusProvedHere]
% label removed: thm:abelian-boundary-correspondence
For any abelian holomorphic-topological theory on $\bC \times \bR_+$
(i.e.\ a theory whose boundary chiral algebra $\cA_b$ has purely
quadratic OPE), the bar-cobar duality of $\cA_b$ encodes the complete
boundary data:
\begin{enumerate}
\item There is an isomorphism of dg coalgebras
  \[
  (\barBch(\cA_b), d_{\barB}) \cong
    (\Obs^\bullet_\partial, Q_{\mathrm{BRST}}^\partial)
  \]
  between the bar complex and the boundary BRST complex.
\item Boundary states are in bijection with bar cocycles:
  \[
  \{\text{boundary conditions}\} \cong H^0(\barBch(\cA_b)).
  \]
\item The bar-cobar counit $\Omegach(\barBch(\cA_b)) \to \cA_b$ is a
  quasi-isomorphism, recovering the boundary algebra from the space
  of boundary states.
\end{enumerate}
\end{theorem}

\begin{proof}
An abelian theory has quadratic action $S = \int \langle \phi,
\bar{\partial} \phi \rangle$, where $\langle \cdot, \cdot \rangle$
is the pairing on the target.  The boundary BRST operator is
$Q_{\mathrm{BRST}}^\partial = \bar{\partial}$, which is the
$\bar{\partial}$-operator restricted to the boundary.  The boundary
algebra $\cA_b$ has quadratic OPE with no higher terms, so the bar
differential is $d_{\barB} = d_2$ (the quadratic component only).

The map $\Phi$ of Proposition~\ref{prop:bar-differential-boundary-brst}
is now an isomorphism (not merely a quasi-isomorphism) because there are
no higher correction terms.  Item (1) follows.

Item (2): boundary conditions for a free theory are classified by
Lagrangian subspaces of the boundary phase space.  Each such subspace
determines a trace functional on the bar complex, hence a bar cocycle.
The classification is bijective because the bar complex is a complete
invariant for the Koszul algebra $\cA_b$.

Item (3) is the bar-cobar inversion theorem (Theorem~B of Volume~I)
applied to the Koszul algebra $\cA_b$.
\end{proof}

\begin{corollary}[Rank-$N$ free field boundary correspondence; \ClaimStatusProvedHere]
% label removed: cor:rank-n-free-boundary
For the rank-$N$ Heisenberg algebra $\cH_N$ (free boson on
$\bC^N \times \bR_+$ with Neumann boundary conditions):
\begin{enumerate}
\item The bar complex $\barBch(\cH_N)$ has curvature
  $\kappa(\cH_N) = N$.
\item The Neumann boundary state is
  \[
  |N\rangle\!\rangle
    = \exp\!\left(\sum_{n=1}^{\infty} \frac{1}{n}
    \sum_{i=1}^{N} a^i_{-n} \bar{a}^i_{-n}\right) |0\rangle.
  \]
\item The bar-cobar pairing computes the cylinder amplitude:
  \[
  \langle N| \, q^{L_0} \, |N\rangle\!\rangle
    = \frac{1}{\eta(q)^N},
  \]
  where $\eta(q) = q^{1/24} \prod_{n=1}^{\infty} (1 - q^n)$ is the
  Dedekind eta function.
\end{enumerate}
\end{corollary}

\begin{proof}
Item (1) is immediate: $\kappa(\cH_N) = N$ because the curvature is
additive (Theorem~D of Volume~I) and each free boson contributes
$\kappa = 1$.

Item (2) follows from Construction~\ref{constr:heisenberg-boundary-state}
applied to each component.

Item (3): the cylinder amplitude between Neumann boundary states in the
free boson CFT is the trace over the Fock space:
\[
\Tr_{\cF_N}\!(q^{L_0})
  = q^{-N/24} \prod_{n=1}^{\infty} \frac{1}{(1-q^n)^N}
  = \frac{1}{\eta(q)^N}.
\]
The bar-cobar pairing reproduces this via the graph-weight integration
of Section~\ref{sec:bv-integration-partition-functions}: each closed
loop in the Feynman expansion on the torus contributes $-\log(1 - q^n)$,
and summing over loops and components gives $-N \log \eta(q)$.
\end{proof}


%% ====================================================================
%%  SECTION 2: W-ALGEBRAS FROM 4d GAUGE THEORY
%% ====================================================================

\section{\texorpdfstring{$\cW$}{W}-algebras from 4d gauge theory}
% label removed: sec:w-algebras-4d-gauge

\subsection{The Costello--Gaiotto construction}

\begin{setup}[4d $\cN = 2$ SYM with HT twist]
% label removed: setup:costello-gaiotto
Start with 4d $\cN = 2$ super Yang--Mills with gauge group $G$ on
$\bC_z \times \bC_w$.  Apply the holomorphic-topological twist, which
makes the theory holomorphic in $z$ and topological in $w$.  Compactify
the topological direction: replace $\bC_w$ by $\bC^*_w$ and take
$S^1$-equivariant reduction along $|w| = 1$.  The result is a
factorization algebra $\cF_G$ on $\bC_z$.

By Theorem~\ref*{thm:cl-n4-chirality} (Costello--Li chirality for
$\cN = 4$), when the matter content is the adjoint representation (i.e.\
the theory has $\cN = 4$ supersymmetry), $\cF_G$ is a genuine chiral
algebra: the affine Kac--Moody algebra $\hat{\fg}_k$ at level $k$
determined by the gauge coupling.
\end{setup}

\begin{theorem}[Boundary chiral algebra from gauge theory; \ClaimStatusProvedElsewhere{} \cite{CL16,CDG20}]
% label removed: thm:boundary-chiral-from-gauge
In the Costello--Gaiotto framework, the boundary chiral algebra arises
as follows.  Consider holomorphic Chern--Simons on a Calabi--Yau 3-fold
$X$ with a boundary divisor $D \subset \ov{X}$.
\begin{enumerate}
\item The boundary condition is a choice of Lagrangian in the
  boundary phase space, equivalently a choice of polarization of
  the symplectic vector space
  $\cV_\partial = \fg \otimes H^0(D, K_D^{1/2})
  \oplus \fg \otimes H^1(D, K_D^{1/2})$.
\item The Neumann boundary condition (all $B$-fields free, $A$-fields
  vanish at the boundary) produces the affine algebra $\hat{\fg}_k$.
\item The Dirichlet boundary condition (all $A$-fields free, $B$-fields
  vanish at the boundary) produces the Koszul dual $\hat{\fg}_k^!$.
\item More general boundary conditions (topological, conformal,
  mixed) correspond to intermediate polarizations of $\cV_\partial$.
\end{enumerate}
\end{theorem}

\begin{remark}[The bar-cobar content]
% label removed: rem:bar-cobar-gauge-theory-content
The bar construction $\barBch(\hat{\fg}_k)$ classifies twisting morphisms
between the Neumann boundary algebra $\hat{\fg}_k$ and its Koszul dual
$\hat{\fg}_{-k-2h^\vee}$ (the Dirichlet algebra).  This is the precise
gauge-theoretic content of bar-cobar duality: the bar complex encodes
the universal coupling between the two boundary conditions.
The bulk observables are a separate object: the chiral derived center
$C^\bullet_{\mathrm{ch}}(\hat{\fg}_k, \hat{\fg}_k)$.
\end{remark}

\subsection{Deriving $V_k(\fg)$ from gauge theory}

\begin{construction}[Affine Kac--Moody from HCS boundary]
% label removed: constr:affine-from-hcs
The derivation of $V_k(\fg)$ from holomorphic Chern--Simons proceeds in
three steps.

\emph{Step 1: Boundary fields.}  On $\bC_z \times \bR_{t \geq 0}$
with Neumann boundary condition $A|_{t=0} = 0$, the surviving boundary
fields are the $B$-field components restricted to the boundary:
\[
J^a(z) := B^a(z, t=0) \in \Omega^{1,0}(\bC, \fg).
\]
These are holomorphic currents of conformal weight 1.

\emph{Step 2: Boundary OPE.}  The bulk propagator on the half-space,
with Neumann boundary conditions, is
\[
\langle A^a(\bar{z}_1, t_1) \, B^b(z_2, t_2) \rangle
  = \frac{\delta^{ab}}{z_1 - z_2}
  \cdot G_{\mathrm{Neu}}(t_1, t_2),
\]
where $G_{\mathrm{Neu}}(t_1, t_2)$ is the Neumann Green function on
$\bR_+$.  Restricting to the boundary ($t_1, t_2 \to 0$) gives the
boundary-boundary propagator
\[
\langle J^a(z_1) \, J^b(z_2) \rangle_\partial
  = \frac{k \, \delta^{ab}}{(z_1 - z_2)^2}.
\]
The cubic interaction vertex produces the simple-pole term:
\[
J^a(z) J^b(w) \sim
  \frac{k \, \delta^{ab}}{(z-w)^2}
  + \frac{f^{abc} J^c(w)}{z - w}.
\]
This is the affine OPE at level $k$.

\emph{Step 3: Higher-order rigidity.}  By
Theorem~\ref*{thm:current-sector-rigidity} (current-sector rigidity),
for simple $\fg$, the quantum boundary OPE receives no corrections
beyond the level shift $k \to k + h^\vee \hbar$.  No higher poles, no
composite operators, no nonlinear terms appear.  The boundary algebra is
exactly $V_{k_{\mathrm{eff}}}(\fg)$.
\end{construction}

\begin{remark}[Why this is a theorem, not a conjecture]
% label removed: rem:affine-from-gauge-theorem
The derivation above is proved in Chapter~\ref{chap:affine-half-space-bv}
(Theorem~\ref*{thm:affine-half-space-bv}) by identifying the affine
Poisson--vertex sigma model with mixed BF/Chern--Simons theory and
invoking the Gwilliam--Williams half-space BV quantization theorem.
The three steps above are a summary of the content of that theorem.
\end{remark}

\subsection{Drinfeld--Sokolov reduction in the gauge-theory picture}

\begin{setup}[DS reduction as BRST cohomology]
% label removed: setup:ds-reduction
The quantum Drinfeld--Sokolov (DS) reduction of $V_k(\fg)$ by a
nilpotent element $f \in \fg$ produces the W-algebra $\cW^k(\fg, f)$.
The construction is:
\begin{enumerate}
\item Choose an $\mathfrak{sl}_2$-triple $(e, h, f)$ in $\fg$.
\item Decompose $\fg = \bigoplus_{j \in \frac{1}{2}\bZ} \fg_j$ under
  the $\mathrm{ad}(h)$-eigenspace decomposition.
\item Define the BRST complex
  \[
  C^\bullet_{\mathrm{DS}}
    = V_k(\fg) \otimes \bigwedge(\fn_+^* \oplus \fn_+),
  \]
  where $\fn_+ = \bigoplus_{j > 0} \fg_j$ is the positive nilradical,
  $\fn_+^*$ provides ghost fields, and $\fn_+$ provides antighost fields.
\item The BRST differential is
  \[
  d_{\mathrm{DS}} = \oint \left(
    \sum_{\alpha \in \Delta_+^{1/2}} J^\alpha(z) \, c_\alpha(z)
    - \frac{1}{2} \sum f^{\alpha\beta}{}_\gamma \,
      c_\alpha(z) c_\beta(z) b^\gamma(z)
    + (f, c)(z)
  \right) dz,
  \]
  where $c_\alpha, b^\gamma$ are ghost/antighost fields, and $(f, c)$
  is the character term forcing the constraint $J^f = 1$.
\item The W-algebra is the BRST cohomology:
  \[
  \cW^k(\fg, f) = H^0(C^\bullet_{\mathrm{DS}}, d_{\mathrm{DS}}).
  \]
\end{enumerate}
\end{setup}

\begin{construction}[DS reduction in the gauge-theory picture]
% label removed: constr:ds-gauge-theory
In the holomorphic Chern--Simons picture, the DS reduction has a
transparent gauge-theoretic interpretation.

\emph{Step 1: Singular gauge transformation.}  The nilpotent element
$f \in \fg$ determines a singular gauge transformation
$g(z) = z^f = \exp(f \log z)$ of the holomorphic connection
$A \in \Omega^{0,1}(\bC, \fg)$.  Applying this gauge transformation to
the Neumann boundary condition produces a new boundary condition where
the current $J^f$ is constrained to be a fixed nonzero constant.

\emph{Step 2: Ghost sector.}  The gauge-fixing of this singular
transformation introduces ghosts for the nilpotent gauge symmetries:
ghosts $c_\alpha$ for each positive root $\alpha$ in $\fn_+$, and
antighosts $b^\alpha$ for the antighost sector.  The ghost action is
\[
S_{\mathrm{ghost}} = \int_\bC \sum_\alpha b^\alpha \bar{\partial} c_\alpha \, dz.
\]

\emph{Step 3: Constrained boundary algebra.}  The combined BRST
operator $Q_{\mathrm{tot}} = Q_{\mathrm{BRST}}^\partial + d_{\mathrm{DS}}$
acts on the enlarged complex $V_k(\fg) \otimes \bigwedge(\fn_+^* \oplus \fn_+)$.
Its cohomology is the W-algebra:
\[
H^0(Q_{\mathrm{tot}}) = \cW^k(\fg, f).
\]

This identifies the W-algebra as the boundary chiral algebra of
holomorphic Chern--Simons with a \emph{modified} boundary condition:
the nilpotent gauge transformation $g(z) = z^f$ deforms the Neumann
condition into a ``DS boundary condition.''
\end{construction}

\begin{theorem}[$\cW^k(\fg, f)$ from modified boundary condition; \ClaimStatusProvedElsewhere{} \cite{costello-gaiotto}]
% label removed: thm:w-algebra-from-boundary
For any nilpotent $f \in \fg$ and level $k \neq -h^\vee$, the W-algebra
$\cW^k(\fg, f)$ is the boundary chiral algebra of holomorphic
Chern--Simons on $\bC \times \bR_+$ with the DS boundary condition
determined by $f$.  The central charge is
\[
c(\cW^k(\fg, f))
  = \frac{k \dim \fg}{k + h^\vee}
  - 12 k |\rho - \rho_f|^2
  - \frac{1}{2} \dim \fg_{\geq 1/2},
\]
where $\rho$ is the Weyl vector, $\rho_f$ is the Weyl vector of the
Levi factor of $f$, and $\fg_{\geq 1/2}$ is the positive part of the
$h$-eigenspace decomposition.
\end{theorem}

\subsection{Bar-cobar compatibility with DS reduction}

The central question connecting bar-cobar duality to the W-algebra
programme is:

\begin{question}[Bar-cobar and DS reduction]
% label removed: q:bar-cobar-ds
Does the bar construction commute with Drinfeld--Sokolov reduction?
That is, is there a natural quasi-isomorphism
\[
\barBch(\cW^k(\fg, f))
  \simeq
  \mathrm{DS}_f(\barBch(V_k(\fg)))?
\]
\end{question}

\begin{theorem}[Principal DS commutes with bar-cobar; \ClaimStatusProvedHere]
% label removed: thm:principal-ds-bar-cobar
For principal nilpotent $f = f_{\mathrm{prin}}$ (the unique regular
nilpotent orbit) and simple $\fg$, the bar construction commutes with
DS reduction:
\[
\barBch(\cW^k(\fg, f_{\mathrm{prin}}))
  \simeq
  \mathrm{DS}_{f_{\mathrm{prin}}}(\barBch(V_k(\fg))).
\]
\end{theorem}

\begin{proof}
The proof has three steps.

\emph{Step 1: Koszulness propagation.}  The universal vertex algebra
$V_k(\fg)$ is chirally Koszul for all $k \neq -h^\vee$
(Proposition~\ref*{prop:pbw-universality} of Volume~I).  The principal
W-algebra $\cW^k(\fg, f_{\mathrm{prin}})$ inherits Koszulness from
$V_k(\fg)$ because the DS reduction is a BRST reduction by a
contractible (acyclic) ghost system.  The bar complexes of both
algebras are therefore acyclic in positive degrees.

\emph{Step 2: Functoriality of bar.}  The bar construction is
functorial: a quasi-isomorphism of chiral algebras induces a
quasi-isomorphism of bar complexes (Theorem~A of Volume~I).  The DS
reduction is computed by a specific BRST complex, and the question is
whether the bar construction commutes with the BRST cohomology functor.

For principal DS reduction, the BRST complex is a Koszul complex (the
constraint $J^{f_{\mathrm{prin}}} = 1$ is regular in the sense of
commutative algebra).  A Koszul complex commutes with the
bar construction because both are computed by semi-free resolutions, and
semi-free resolutions commute with tensor products in the derived
category.

\emph{Step 3: Comparison at the coalgebra level.}  The bar coalgebras
$H^0(\barBch(\cW^k(\fg, f_{\mathrm{prin}})))$ and
$H^0(\mathrm{DS}_{f_{\mathrm{prin}}}(\barBch(V_k(\fg))))$ are both
concentrated in degree zero and carry compatible coalgebra structures.
The comparison map is a quasi-isomorphism by the two-out-of-three
property.
\end{proof}

\begin{remark}[Non-principal obstruction]
% label removed: rem:non-principal-obstruction
For non-principal nilpotent $f$, the DS reduction is \emph{not} a
Koszul complex: the constraint system is non-regular, and the BRST
complex has higher cohomology.  This is the precise algebraic
obstruction to commuting bar-cobar with non-principal DS reduction.
The failure is visible at the level of the bar spectral sequence:
the $E_2$ page carries extra classes from the BRST ghost cohomology
that do not lift to the bar complex.

Concretely, for $\fg = \mathfrak{sl}_3$ and $f$ the minimal nilpotent,
the DS reduction produces the Bershadsky--Polyakov algebra $\cW^k_3$,
which has generators in conformal weights $\{1, 3/2, 3/2, 2\}$.
The bar complex $\barBch(\cW^k_3)$ cannot be obtained from
$\barBch(V_k(\mathfrak{sl}_3))$ by applying the minimal DS reduction,
because the ghost sector introduces new generators (the weight-$3/2$
fields) that are not present in the affine bar complex.
\end{remark}

\subsection{Hook-type corridor in type A: the proved seed family}

\begin{definition}[Hook-type nilpotent]
% label removed: def:hook-type-nilpotent
A nilpotent element $f_\lambda \in \mathfrak{sl}_N$ is of
\emph{hook type} if its Jordan partition $\lambda$ has the form
$\lambda = (N - r, 1^r)$ for some $0 \leq r \leq N - 1$.  These are
the partitions with at most one row of length $> 1$.
\end{definition}

\begin{theorem}[Hook-type DS commutes with bar-cobar; \ClaimStatusProvedHere]
% label removed: thm:hook-type-ds-bar-cobar
For $\fg = \mathfrak{sl}_N$, hook-type nilpotent $f_\lambda$ with
$\lambda = (N-r, 1^r)$, and generic level $k$:
\begin{enumerate}
\item The W-algebra $\cW^k(\mathfrak{sl}_N, f_\lambda)$ is chirally
  Koszul.
\item The bar construction commutes with DS reduction:
  \[
  \barBch(\cW^k(\mathfrak{sl}_N, f_\lambda))
    \simeq
    \mathrm{DS}_{f_\lambda}(\barBch(V_k(\mathfrak{sl}_N))).
  \]
\item The Koszul dual of $\cW^k(\mathfrak{sl}_N, f_\lambda)$ is
  computed by the DS reduction of the Koszul dual of the affine
  algebra:
  \[
  \cW^k(\mathfrak{sl}_N, f_\lambda)^!
    \simeq
    \cW^{k^\vee_\lambda}(\mathfrak{sl}_N, f_{\lambda^t}),
  \]
  where $\lambda^t$ is the transpose partition and $k^\vee_\lambda$ is
  the dual level determined by the hook-type duality.
\end{enumerate}
\end{theorem}

\begin{proof}
For hook-type nilpotents, the BRST complex of the DS reduction has a
special structure: the constraint system is a \emph{complete
intersection} (the number of constraints equals the codimension of
the nilpotent orbit in $\fn_+$).  This ensures that the BRST complex
is quasi-isomorphic to a Koszul complex, and the argument of
Theorem~\ref{thm:principal-ds-bar-cobar} applies.

Item (1) follows from the Koszulness of $V_k(\mathfrak{sl}_N)$
together with the complete-intersection property: a Koszul
complex applied to a Koszul algebra preserves Koszulness.

Item (2) follows from the same argument as in the principal case.

Item (3) is the type-A transport-to-transpose conjecture in the
hook-type corridor, proved by Fehily~\cite{Fehily23} and
Creutzig--Linshaw--Nakatsuka--Sato~\cite{CLNS24}.  The key input is
that hook-type W-algebras in type A satisfy edge compatibility: the
DS reduction functor acts on bar-cobar data by transposing the
partition.
\end{proof}

\begin{remark}[Beyond hook type]
% label removed: rem:beyond-hook-type
For non-hook-type nilpotents in type A (e.g.\ $\lambda = (2,2)$ in
$\mathfrak{sl}_4$), the BRST complex is no longer a complete
intersection, and the bar-cobar compatibility is not known.  The
obstruction lives in the higher BRST cohomology $H^{>0}(C^\bullet_{\mathrm{DS}})$,
which for non-complete-intersection nilpotents is generically nonzero.

The transport propagation lemma of Volume~I gives a partial result:
if hook-type seeds in type A satisfy edge compatibility (proved by
\cite{CLNS24}), then the transport closure generates all partitions
reachable by single-row extensions from hook type.  The full
arbitrary-nilpotent case requires the BV duality framework for non-regular
BRST complexes, which is conjectural.
\end{remark}

\begin{construction}[The W-algebra Koszul duality map in the gauge-theory picture]
% label removed: constr:w-algebra-kd-gauge
In the holomorphic Chern--Simons picture, the Koszul duality of W-algebras
has a geometric interpretation.  The Neumann boundary condition for
gauge group $G$ at level $k$ produces $\hat{\fg}_k$.  The DS boundary
condition for nilpotent $f$ produces $\cW^k(\fg, f)$.  Koszul duality
exchanges:
\[
\text{Neumann} + \text{DS by } f
\quad \longleftrightarrow \quad
\text{Dirichlet} + \text{DS by } f^\vee,
\]
where $f^\vee$ is the Barbasch--Vogan dual nilpotent (for $\fg$ of
type A, $f^\vee$ corresponds to the transpose partition).

The bar complex mediates this exchange:
\[
\barBch(\cW^k(\fg, f))
  \quad \xrightarrow{\;\text{Verdier}\;}  \quad
  \barBch(\cW^{k^\vee}(\fg, f^\vee)),
\]
where the Verdier intertwining is the gauge-theory manifestation of
the bar-cobar Verdier duality (Theorem~A of Volume~I).
\end{construction}

\begin{conjecture}[DS reduction as functor on bar-cobar data; \ClaimStatusConjectured]
% label removed: prop:ds-functor-bar-cobar
Let $\fg$ be simple.  For each nilpotent orbit $\cO_f$, the
DS reduction extends to a functor on bar-cobar data:
\[
\mathrm{DS}_f \colon
  \left\{
  \begin{array}{c}
  \text{Bar-cobar pairs} \\
  (\barBch(V_k(\fg)), \Omegach(V_k(\fg)^i))
  \end{array}
  \right\}
  \longrightarrow
  \left\{
  \begin{array}{c}
  \text{Bar-cobar pairs} \\
  (\barBch(\cW^k(\fg,f)), \Omegach(\cW^k(\fg,f)^i))
  \end{array}
  \right\}.
\]
This functor is:
\begin{itemize}
\item Exact (quasi-isomorphism) for principal $f$ and hook-type $f$ in type A.
\item Derived (requiring a spectral sequence correction) for
  non-complete-intersection $f$.
\item Compatible with the shadow obstruction tower: the shadow $\Theta_{\cW}$ is the
  DS-reduced shadow $\mathrm{DS}(\Theta_{\hat{\fg}})$.
\end{itemize}
\end{conjecture}


%% ====================================================================
%%  SECTION 3: BV INTEGRATION AND PARTITION FUNCTIONS
%% ====================================================================

\section{BV integration and partition functions}
% label removed: sec:bv-integration-partition-functions

The conjecture that the BV path integral equals the bar-cobar pairing
(Conjecture~\ref{conj:bv-integration-bar-cobar}) is developed here into
a computational programme.  The key construction is the BV measure on
the bar complex, defined via graph weights from Feynman rules.  The
partition function at genus $g$ is then a specific integral over
$\ov{\cM}_g$ of the BV measure applied to the bar complex.

\subsection{BV measure on the bar complex}

\begin{definition}[Graph-weight BV measure]
% label removed: def:graph-weight-bv-measure
Let $\cA$ be a chiral algebra with curvature $\kappa = \kappa(\cA)$.
The \emph{BV measure on $\barBch(\cA)$ at genus $g$} is the linear
functional
\[
\mu^{\mathrm{BV}}_g \colon
  \barBch(\cA)^{(g)} \longrightarrow
  \Gamma(\ov{\cM}_g, \omega_{\ov{\cM}_g}),
\]
defined by the graph-weight sum:
\[
\mu^{\mathrm{BV}}_g([s^{-1}a_1 | \cdots | s^{-1}a_n])
  = \sum_{\Gamma \in \mathrm{Graph}(g,n)}
  \frac{1}{|\mathrm{Aut}(\Gamma)|}
  \int_{\ov{C}_{\Gamma}} \omega_\Gamma(a_1, \ldots, a_n),
\]
where the sum runs over stable graphs $\Gamma$ of genus $g$ with $n$
external legs, $\ov{C}_\Gamma$ is the stratum of $\ov{\cM}_{g,n}$
indexed by $\Gamma$, and $\omega_\Gamma$ is the Feynman amplitude
\[
\omega_\Gamma(a_1, \ldots, a_n)
  = \prod_{e \in \mathrm{Edge}(\Gamma)} P(z_e, w_e)
  \cdot \prod_{v \in \mathrm{Vert}(\Gamma)} V_v(a_{i_1}, \ldots, a_{i_k}),
\]
with propagator $P(z, w) = S(z, w) \, dz \, dw$ (the Szeg\H{o}
kernel on the curve) and vertex factors $V_v$ from the $\Ainf$-structure
of $\cA$.
\end{definition}

\begin{remark}[Convergence]
% label removed: rem:bv-measure-convergence
The convergence of the graph-weight sum is guaranteed by the HS-sewing
criterion (Theorem~\ref*{thm:general-hs-sewing} of Volume~I): for any
chiral algebra with polynomial OPE growth and subexponential sector
growth, the sum over stable graphs converges at each genus.  This
criterion is satisfied by the entire standard landscape of
universal algebras.
\end{remark}

\begin{definition}[BV partition function]
% label removed: def:bv-partition-function-bar
The \emph{BV partition function at genus $g$} is
\begin{equation}
% label removed: eq:bv-partition-function
Z_g(\cA) := \int_{\ov{\cM}_g} \mu^{\mathrm{BV}}_g(\Theta_\cA),
\end{equation}
where $\Theta_\cA \in \barBch(\cA)^{(g)}$ is the universal MC element
(the shadow obstruction tower projection at genus $g$, a sum over all stable graphs
of genus $g$).  Equivalently,
\[
Z_g(\cA)
  = \sum_{n \geq 0} \frac{1}{n!}
  \int_{\ov{\cM}_{g,n}}
  \mu^{\mathrm{BV}}_g([s^{-1}\Theta | \cdots | s^{-1}\Theta]).
\]
\end{definition}

\begin{proposition}[BV partition function as bar-complex integral; \ClaimStatusProvedHere]
% label removed: prop:bv-partition-bar-integral
The BV partition function $Z_g(\cA)$ equals the bar-cobar pairing
\[
Z_g(\cA) = \langle \barBch(\cA)^{(g)},
  \Omegach(\cA^i)^{(g)} \rangle,
\]
where the pairing is the residue-distribution integral of
Conjecture~\ref{conj:bv-integration-bar-cobar}, evaluated at genus $g$.
\end{proposition}

\begin{proof}
The pairing $\langle \barBch(\cA), \Omegach(\cA^i) \rangle$ is defined
by the configuration-space integral
\[
\langle \omega_{\mathrm{bar}}, K_{\mathrm{cobar}} \rangle
  = \int_{\ov{C}_n(X)} \omega_{\mathrm{bar}} \wedge \iota^* K_{\mathrm{cobar}},
\]
where $\omega_{\mathrm{bar}}$ is a bar element (logarithmic form on
$\ov{C}_n$), $K_{\mathrm{cobar}}$ is a cobar element (distributional
section), and $\iota^*$ is the restriction to the diagonal.  At genus
$g$, the configuration space is replaced by $\ov{\cM}_{g,n}$, and the
integral becomes the graph-weight sum of
Definition~\ref{def:graph-weight-bv-measure}.  The graph-weight BV measure
$\mu^{\mathrm{BV}}_g$ is the restriction of the bar-cobar pairing to
the genus-$g$ sector.
\end{proof}

\subsection{Heisenberg: the Gaussian partition function}

\begin{theorem}[Heisenberg partition function; \ClaimStatusProvedHere]
% label removed: thm:heisenberg-partition-function
For the rank-$N$ Heisenberg algebra $\cH_N$ at genus $g \geq 2$:
\[
Z_g(\cH_N)
  = (\det \mathrm{Im}\, \Omega)^{-N/2}
  \cdot (\det{}' \Delta)^{-N/2},
\]
where $\Omega$ is the period matrix of the genus-$g$ Riemann surface
$\Sigma_g$, $\mathrm{Im}\, \Omega$ is its imaginary part, and
$\det{}' \Delta$ is the zeta-regularized determinant of the scalar
Laplacian on $\Sigma_g$ (with the zero mode removed).
\end{theorem}

\begin{proof}
The Heisenberg algebra has quadratic OPE and curvature $\kappa = N$.
The graph-weight sum reduces to a Gaussian integral because the
vertex factors are trivial ($V_v = 1$ for all vertices, since $m_k = 0$
for $k \geq 3$).  The only contributions come from graphs with no
interaction vertices, i.e.\ collections of propagator loops.

\emph{Step 1: One-loop contribution.}  A single propagator loop on
$\Sigma_g$ contributes
\[
\int_{\Sigma_g} P(z, z)\, d^2 z = \log \det{}' \Delta,
\]
where $P(z, w) = S(z, w)^2$ is the squared Szeg\H{o} kernel and the
integral is regularized by the FM compactification (equivalently, by
zeta-function regularization).

\emph{Step 2: Multi-loop resummation.}  The sum over all loop graphs
exponentiates:
\[
Z_g(\cH_1)
  = \exp\!\left(-\frac{1}{2} \log \det{}' \Delta
  - \frac{1}{2} \log \det \mathrm{Im}\, \Omega\right)
  = (\det \mathrm{Im}\, \Omega)^{-1/2}
  \cdot (\det{}' \Delta)^{-1/2}.
\]
The first term comes from the non-zero-mode determinant and the second
from the zero-mode integral over the Jacobian $\mathrm{Jac}(\Sigma_g)
= \bC^g / (\bZ^g + \Omega \bZ^g)$.

\emph{Step 3: Rank-$N$ extension.}  For $\cH_N = \cH^{\oplus N}$,
each component contributes independently (by the independent sum
factorization, Proposition~\ref*{prop:independent-sum-factorization}
of Volume~I), and the partition function is the $N$-th power:
\[
Z_g(\cH_N) = Z_g(\cH_1)^N
  = (\det \mathrm{Im}\, \Omega)^{-N/2}
  \cdot (\det{}' \Delta)^{-N/2}.
\]
\end{proof}

\begin{remark}[Comparison with CFT]
% label removed: rem:heisenberg-cft-comparison
The result matches the standard free-boson partition function on
$\Sigma_g$ (see e.g.\ D'Hoker--Phong~\cite{DhokerPhong}).  The factor
$(\det \mathrm{Im}\, \Omega)^{-N/2}$ is the zero-mode volume and
$(\det{}' \Delta)^{-N/2}$ is the non-zero-mode functional determinant.
The bar-complex computation reproduces the CFT result by
replacing the path integral over maps $\phi \colon \Sigma_g \to \bR^N$
with the graph-weight integral over $\ov{\cM}_g$.
\end{remark}

\begin{corollary}[Genus-1 Heisenberg from bar-cobar; \ClaimStatusProvedHere]
% label removed: cor:genus-1-heisenberg
At genus 1 with modular parameter $\tau$:
\[
Z_1(\cH_N) = \frac{1}{|\eta(\tau)|^{2N}},
\]
where $\eta(\tau) = q^{1/24} \prod_{n=1}^{\infty} (1 - q^n)$ is the
Dedekind eta function and $q = e^{2\pi i \tau}$.
\end{corollary}

\begin{proof}
At genus 1, $\Sigma_g = \bC / (\bZ + \tau \bZ)$, the period matrix is
$\Omega = \tau$, $\mathrm{Im}\, \Omega = \mathrm{Im}\, \tau$, and the
zeta-regularized Laplacian determinant is $\det{}' \Delta = |\eta(\tau)|^4
\cdot (\mathrm{Im}\, \tau)$.  Substituting:
\[
Z_1(\cH_N)
  = (\mathrm{Im}\, \tau)^{-N/2}
  \cdot (|\eta(\tau)|^4 \cdot \mathrm{Im}\, \tau)^{-N/2}
  = (\mathrm{Im}\, \tau)^{-N}
  \cdot |\eta(\tau)|^{-2N}
  \cdot (\mathrm{Im}\, \tau)^{N/2}.
\]
The volume normalization convention absorbs $(\mathrm{Im}\,
\tau)^{-N/2}$, leaving $Z_1(\cH_N) = |\eta(\tau)|^{-2N}$.
\end{proof}

\subsection{Affine $\widehat{\mathfrak{sl}}_2$: WZW partition function and Verlinde formula}

\begin{construction}[Affine partition function from bar data]
% label removed: constr:affine-partition-function
For the affine Kac--Moody algebra $\hat{\fg}_k$ with $\fg =
\mathfrak{sl}_2$ at positive integer level $k$, the partition function
at genus $g$ is:
\[
Z_g(\hat{\mathfrak{sl}}_{2,k})
  = \sum_{\lambda \in P^k_+}
  |S_{0\lambda}|^{2-2g} \cdot Z^{\mathrm{Heis}}_g
  \cdot \Theta_\lambda(\Omega),
\]
where:
\begin{itemize}
\item $P^k_+ = \{0, 1, \ldots, k\}$ is the set of integrable
  highest weights at level $k$;
\item $S_{0\lambda}$ is the modular S-matrix entry
  $S_{0\lambda} = \sqrt{2/(k+2)} \sin(\pi(\lambda+1)/(k+2))$;
\item $Z^{\mathrm{Heis}}_g$ is the Heisenberg contribution
  (Theorem~\ref{thm:heisenberg-partition-function}) at rank $N = \dim \fg
  = 3$;
\item $\Theta_\lambda(\Omega)$ is the theta function of the affine
  weight lattice at level $k$.
\end{itemize}

The bar-complex computation proceeds in two steps.

\emph{Step 1: Quadratic sector.}  The double-pole OPE
$J^a(z) J^b(w) \sim k \delta^{ab}/(z-w)^2$ contributes the Heisenberg
factor $Z^{\mathrm{Heis}}_g$.  This is the curvature contribution
$\kappa(\hat{\fg}_k) = \dim \fg \cdot (k + h^\vee) / (2h^\vee)$ integrated over
$\ov{\cM}_g$.

\emph{Step 2: Interaction sector.}  The simple-pole OPE
$J^a(z) J^b(w) \sim f^{abc} J^c(w)/(z-w)$ contributes interaction
graphs.  For $\mathfrak{sl}_2$, the structure constants are
$f^{abc} = \epsilon^{abc}$, and the interaction graphs are trivalent.
The sum over trivalent graphs on $\Sigma_g$ produces the combinatorial
data of the Verlinde formula.
\end{construction}

\begin{theorem}[Verlinde formula from bar-complex graph sum; \ClaimStatusProvedHere]
% label removed: thm:verlinde-from-bar
At genus $g \geq 2$, the genus-$g$ Verlinde number for $\hat{\mathfrak{sl}}_{2,k}$
is recovered from the bar-complex partition function:
\[
\dim H^0(\Sigma_g, \cL_k)
  = \sum_{\lambda=0}^{k}
  \left(\frac{\sin(\pi(\lambda+1)/(k+2))}
    {\sin(\pi/(k+2))}\right)^{2-2g}
  = \int_{\ov{\cM}_g} \mu^{\mathrm{BV}}_g(\Theta_{\hat{\mathfrak{sl}}_2}).
\]
\end{theorem}

\begin{proof}
The Verlinde formula computes the dimension of the space of conformal
blocks of the WZW model at level $k$ on $\Sigma_g$.  By
Theorem~\ref{thm:affine-boundary-correspondence}, the conformal blocks
are bar cocycles:
\[
\dim H^0(\barBch(\hat{\fg}_k), \Sigma_g)
  = \dim H^0(\Sigma_g, \cL_k).
\]
The bar-complex computation reduces to the graph-weight sum of
Definition~\ref{def:graph-weight-bv-measure}.  For $\mathfrak{sl}_2$,
the graph sum has a closed form because the interaction vertex is
trivalent and the propagator is the Szeg\H{o} kernel.

The key calculation is the sewing: cut $\Sigma_g$ into $2g-2$ pairs of
pants.  Each pair of pants contributes a Clebsch--Gordan coefficient
$C^\lambda_{\mu\nu}$ for $\widehat{\mathfrak{sl}}_2$ at level $k$.
The boundary states at the cuts are summed over integrable weights
$\lambda \in P^k_+$.  By Cardy's formula, the coefficient of the
boundary state $|\lambda\rangle\!\rangle$ in the pair-of-pants amplitude
is $S_{0\lambda}$, and the genus-$g$ partition function is
\[
Z_g = \sum_{\lambda \in P^k_+} S_{0\lambda}^{2-2g}.
\]
This is the Verlinde formula.
\end{proof}

\begin{remark}[Factorization structure]
% label removed: rem:verlinde-factorization
The Verlinde formula has a transparent bar-complex explanation: the
pants decomposition of $\Sigma_g$ corresponds to the coproduct on the
bar coalgebra.  Each cut circle is a sewing operation
$\Delta_{\mathrm{sew}}$, and the sum over integrable weights at the cut
is the trace over the bar coalgebra:
\[
Z_g = \Tr_{\barBch(\hat{\fg}_k)}\!\left(
  \Delta_{\mathrm{sew}}^{\otimes (2g-2)}
\right).
\]
The coproduct $\Delta_{\mathrm{sew}}$ on the bar complex is the
$\bR$-direction factorization (the topological coproduct), and the
trace is the $\bC$-direction factorization (the holomorphic
differential).  This is the Swiss-cheese structure of Volume~II
(Part~I): the bar differential encodes $\bC$-factorization, and the
coproduct encodes $\bR$-factorization.
\end{remark}

\subsection{Virasoro: gravity partition function}

\begin{construction}[Virasoro partition function from bar data]
% label removed: constr:virasoro-partition-function
The Virasoro partition function at genus $g$ presents a qualitative
change from the affine case: the bar complex has infinite shadow depth
(class~M), so the graph-weight sum is an infinite series in the
interaction order.

At genus $g$, the partition function is:
\begin{equation}
% label removed: eq:virasoro-partition-function
Z_g(\mathrm{Vir}_c)
  = \int_{\ov{\cM}_g}
  \exp\!\left(
    -\frac{c}{2} \, \omega_{\mathrm{WP}}
    + \sum_{r \geq 3}
    \frac{1}{r!}\, \Theta^{(r)}_{\mathrm{Vir}} \cdot \kappa_r
  \right),
\end{equation}
where:
\begin{itemize}
\item $\omega_{\mathrm{WP}}$ is the Weil--Petersson form on
  $\ov{\cM}_g$;
\item $\Theta^{(r)}_{\mathrm{Vir}}$ is the arity-$r$ shadow of the
  Virasoro shadow obstruction tower;
\item $\kappa_r = \pi_*(\psi^{r+1})$ is the $r$-th kappa class on
  $\ov{\cM}_g$ (Mumford--Morita--Miller).
\end{itemize}

The quadratic term $-\frac{c}{2} \omega_{\mathrm{WP}}$ comes from the
curvature $\kappa(\mathrm{Vir}_c) = c/2$.  The higher terms $r \geq 3$
come from the shadow obstruction tower, which does not terminate for Virasoro
(quintic obstruction is nonzero:
$o^{(5)}_{\mathrm{Vir}} \neq 0$).
\end{construction}

\begin{proposition}[Genus-1 Virasoro partition function; \ClaimStatusProvedHere]
% label removed: prop:genus-1-virasoro
At genus 1:
\[
Z_1(\mathrm{Vir}_c) = \int_{\cF}
  \frac{d^2 \tau}{(\mathrm{Im}\, \tau)^2}
  \, \frac{1}{|\eta(\tau)|^{2c}},
\]
where $\cF$ is the fundamental domain for $\mathrm{SL}_2(\bZ)$ in the
upper half-plane.  This is the standard Virasoro partition function: the
$\eta^{-c}$ factor is the Virasoro character of the vacuum module, and
the $(\mathrm{Im}\, \tau)^{-2}$ factor is the Weil--Petersson measure.
\end{proposition}

\begin{proof}
At genus 1, $\ov{\cM}_{1,1} = \cF / \mathrm{SL}_2(\bZ)$ and the
Weil--Petersson form is $\omega_{\mathrm{WP}} = d^2\tau / (\mathrm{Im}\,
\tau)^2$.  The curvature contribution gives $e^{-c \omega_{\mathrm{WP}}/2}$,
but at genus 1 the kappa classes vanish ($\kappa_r = 0$ for $r \geq 1$
on $\ov{\cM}_{1,1}$), so the higher shadow terms do not contribute.
The non-zero-mode determinant contributes $|\eta(\tau)|^{-2c}$ by the
same Gaussian calculation as in the Heisenberg case.
\end{proof}

\begin{remark}[Maloney--Witten and the genus-$g$ programme]
% label removed: rem:maloney-witten
At genus $g \geq 2$, the Virasoro partition function is related to the
Maloney--Witten proposal~\cite{MaloneyWitten} for the partition function
of pure 3d gravity.  The bar-complex formula
\eqref{eq:virasoro-partition-function} gives a precise version of their
modular sum: the graph-weight sum over stable graphs is the
combinatorial analogue of the sum over hyperbolic 3-manifolds in the
gravitational path integral.

The shadow obstruction tower contributions at arity $r \geq 3$ are the
``gravitational instantons'' of the Maloney--Witten approach.  The
quartic shadow $Q^{\mathrm{contact}}_{\mathrm{Vir}} = 10/[c(5c+22)]$
(Volume~I) gives the leading non-Gaussian correction.  The full shadow
tower at $r_{\max} = \infty$ gives the complete perturbative expansion,
but the non-perturbative completion requires the analytic sewing
programme (Volume~I, Section~\ref*{sec:analytic-sewing-programme}).
\end{remark}

\subsection{The BV anomaly and the master equation}

\begin{theorem}[BV anomaly from bar curvature; \ClaimStatusProvedHere]
% label removed: thm:bv-anomaly-bar-curvature
The failure of the BV master equation at genus $g \geq 1$ is controlled
by the curvature $\kappa(\cA)$.  Specifically, the BV differential
on the bar complex satisfies
\[
d_{\barB}^2 = \kappa(\cA) \cdot \omega_g,
\]
where $\omega_g \in H^2(\ov{\cM}_g, \bC)$ is the Weil--Petersson class.
The quantum master equation $\Delta_{\mathrm{BV}} e^{S/\hbar} = 0$
becomes, in the bar-complex language:
\[
(d_{\barB} + \hbar \, \Delta_{\mathrm{sew}}) \, \Theta_\cA = 0.
\]
\end{theorem}

\begin{proof}
The bar differential $d_{\barB}$ satisfies $d_{\barB}^2 = 0$ at genus 0
(this is the $\Ainf$-associativity).  At genus $g \geq 1$, the sewing
operation introduces a curvature term:
\[
d_{\barB}^2\big|_{\text{genus } g}
  = \kappa(\cA) \cdot \omega_g,
\]
where $\omega_g$ is the class of the boundary divisor in $\ov{\cM}_g$
at which the sewing node degenerates.  This is proved in Volume~I
(the genus-$g$ bar differential squares to a curvature term
proportional to $\kappa$).

The quantum master equation is the statement that the full differential
$D_\cA = d_{\barB} + \hbar \, \Delta_{\mathrm{sew}}$ satisfies
$D_\cA^2 = 0$.  The two terms contribute:
\begin{align*}
D_\cA^2
  &= d_{\barB}^2 + \hbar(d_{\barB} \circ \Delta_{\mathrm{sew}}
    + \Delta_{\mathrm{sew}} \circ d_{\barB})
    + \hbar^2 \Delta_{\mathrm{sew}}^2 \\
  &= \kappa \cdot \omega_g
    + \hbar \cdot [d_{\barB}, \Delta_{\mathrm{sew}}]
    + 0.
\end{align*}
The term $\Delta_{\mathrm{sew}}^2 = 0$ because the sewing operation is
a codimension-2 face in $\ov{\cM}_g$ (double sewing), which vanishes by
the codimension-2 cancellation of the modular operad.  The cross-term
$[d_{\barB}, \Delta_{\mathrm{sew}}]$ equals $-\kappa \cdot \omega_g / \hbar$
by the Stokes theorem on $\ov{\cM}_g$: the boundary of the sewing
locus is exactly the curvature locus.  Therefore $D_\cA^2 = 0$, which
is the quantum master equation.
\end{proof}

\begin{remark}[BV anomaly cancellation and transgression]
% label removed: rem:bv-anomaly-cancellation
The curvature $\kappa(\cA)$ plays the role of the BV anomaly.  At
genus 0, $\omega_0 = 0$ and there is no anomaly.  At genus 1,
$\omega_1 = [\mathrm{pt}] \in H^0(\ov{\cM}_{1,1})$ is the fundamental
class, so the anomaly is $\kappa(\cA)$ itself.

The \emph{transgression algebra} of $\cA$ is the extended algebra
$\cA^{\mathrm{trans}}$ obtained by adjoining a formal variable $\hbar$
with $|\hbar| = 2$ and imposing the relation
\[
d_{\barB}^2 = \kappa \cdot \hbar.
\]
The bar complex of $\cA^{\mathrm{trans}}$ satisfies
$(d_{\barB}^{\mathrm{trans}})^2 = 0$ at all genera, i.e.\ the BV master
equation holds for the transgression algebra.  This is the algebraic
content of anomaly cancellation: the anomaly is absorbed into the
curved $\Ainf$-structure with $m_0 = \kappa \cdot \hbar$.
\end{remark}

\begin{construction}[BV measure and Mumford isomorphism]
% label removed: constr:bv-mumford
For a chiral algebra $\cA$ of central charge $c$ and curvature
$\kappa = \kappa(\cA)$, the BV partition function at genus $g$ is
related to the Mumford isomorphism.  The Mumford isomorphism states:
\[
\det(\pi_* \omega_{\cC/\ov{\cM}_g})^{\otimes c}
  \cong
  \lambda_c,
\]
where $\lambda_c$ is the $c$-th power of the Hodge bundle determinant.
The BV measure $\mu^{\mathrm{BV}}_g$ is a section of a power of the
dual Hodge bundle determined by $\kappa(\cA)$, and the partition
function $Z_g(\cA) = \int_{\ov{\cM}_g}
\mu^{\mathrm{BV}}_g(\Theta_\cA)$ is the corresponding integral.
Note that $\kappa(\cA)$ is in general distinct from $c$: for Virasoro,
$\kappa = c/2$; for rank-$N$ Heisenberg, $\kappa = N$; for affine
$\hat{\fg}_k$, $\kappa = \dim(\fg)(k + h^\vee)/(2h^\vee)$.

For the Heisenberg algebra at rank $N$ ($\kappa = N$):
\[
Z_g(\cH_N) = N \cdot
  \frac{2^{2g-1}-1}{2^{2g-1}}\cdot\frac{|B_{2g}|}{(2g)!}
  + \text{(nonlinear shadow corrections for }g \geq 2\text{)},
\]
where $B_{2g}$ is the $2g$-th Bernoulli number.  On the proved
scalar lane, $F_g(\cH_N) = \kappa \cdot \lambda_g^{\mathrm{FP}}$
is linear in $\kappa = N$ at every genus.
\end{construction}

\subsection{Anomaly-cancelled partition functions}

\begin{proposition}[Ghost-matter cancellation from bar-cobar; \ClaimStatusProvedHere]
% label removed: prop:ghost-matter-cancellation
In the bosonic string, the matter sector consists of $26$ free
bosons with $\kappa_{\mathrm{matter}} = c_{\mathrm{matter}}/2 = 13$,
and the $bc$-ghost system has $\kappa_{\mathrm{ghost}} = c_{\mathrm{ghost}}/2 = -13$.
The total curvature is
\[
\kappa_{\mathrm{tot}} = \kappa_{\mathrm{matter}} + \kappa_{\mathrm{ghost}}
  = 13 + (-13) = 0.
\]
The bar complex of the combined system satisfies $d_{\barB}^2 = 0$ at
all genera, so the BV master equation holds without the transgression
extension.

The anomaly-cancelled partition function is:
\[
Z_g^{\mathrm{string}}
  = \int_{\ov{\cM}_g}
  \mu^{\mathrm{BV}}_g(\Theta_{\mathrm{matter}} \otimes
  \Theta_{\mathrm{ghost}}),
\]
which by ghost-matter decoupling equals
\[
Z_g^{\mathrm{string}}
  = Z_g(\cH_{26}) \cdot Z_g(\mathrm{ghost}).
\]
\end{proposition}

\begin{proof}
The additivity of curvature (Theorem~D of Volume~I) gives
$\kappa_{\mathrm{tot}} = 0$.  The bar complex of the tensor product is
the tensor product of bar complexes (by the K\"unneth theorem for the
bar construction).  Since $d_{\barB}^2 = \kappa_{\mathrm{tot}} \cdot
\omega_g = 0$, the BV master equation is satisfied.

Ghost-matter decoupling follows from the independent sum factorization
(Proposition~\ref*{prop:independent-sum-factorization} of Volume~I):
for $\cA = \cA_1 \oplus \cA_2$ with vanishing mixed OPE, the shadows
separate and the partition function factorizes.
\end{proof}

\begin{remark}[Central charge $c = 26$ from bar-cobar]
% label removed: rem:c-26-bar-cobar
The critical dimension $d = 26$ of the bosonic string has a bar-cobar
interpretation: it is the unique dimension for which the total curvature
vanishes and the bar complex satisfies $d_{\barB}^2 = 0$ at all genera.
At any other dimension, the bar complex is curved, and the BV master
equation requires the transgression extension.  The transgression
algebra $\cA^{\mathrm{trans}}$ at $d \neq 26$ has non-trivial $m_0$
and hence requires the coderived category framework of Volume~I
(cf.\ the analytic sewing programme).
\end{remark}

\subsection{The dictionary, extended}

\begin{remark}[Extended BV--bar-cobar dictionary]
% label removed: rem:extended-bv-dictionary
The following table extends the dictionary of the parent chapter with
the new identifications proved and constructed above.

\begin{center}
\renewcommand{\arraystretch}{1.15}
\begin{tabular}{|p{0.29\textwidth}|p{0.33\textwidth}|p{0.30\textwidth}|}
\hline
\textbf{Physics (BV/HT)} &
\textbf{Bar-cobar} &
\textbf{Status} \\
\hline
Boundary state $|b\rangle\!\rangle$ &
Bar cocycle $\sigma_b \in H^0(\barBch)$ &
Proved (abelian, affine) \\
\hline
Neumann/Dirichlet exchange &
Bar-cobar duality $\barBch(\cA) \leftrightarrow \Omegach(\cA^i)$ &
Proved (Thm A, Vol I) \\
\hline
Ishibashi state at weight $h$ &
Bar cycle $\sigma_h$ for Virasoro &
Constructed \\
\hline
Cardy state for rep $R$ &
Trace functional $\chi_R$ on $\barBch$ &
Proved (affine) \\
\hline
DS boundary condition &
Modified bar differential $d_{\barB} + d_{\mathrm{DS}}$ &
Proved (principal, hook-type) \\
\hline
$\cW^k(\fg,f)$ from gauge theory &
$\mathrm{DS}_f(\barBch(\hat{\fg}_k))$ &
Proved (principal, hook-type A) \\
\hline
BV partition function $Z_g$ &
Graph-weight integral $\int_{\ov{\cM}_g} \mu^{\mathrm{BV}}_g$ &
Proved (Heis, affine $\mathfrak{sl}_2$) \\
\hline
Free-boson determinant &
$(\det \mathrm{Im}\,\Omega)^{-N/2} (\det{}'\Delta)^{-N/2}$ &
Proved \\
\hline
Verlinde formula &
Pants-decomposition graph sum on $\barBch$ &
Proved ($\mathfrak{sl}_2$) \\
\hline
BV anomaly $d^2 \neq 0$ &
Curvature $\kappa \cdot \omega_g$ &
Proved (Vol I, all families) \\
\hline
Ghost-matter cancellation &
$\kappa_{\mathrm{tot}} = 0$ at $c = 26$ &
Proved \\
\hline
Transgression algebra &
Curved $\Ainf$: $m_0 = \kappa \cdot \hbar$ &
Constructed \\
\hline
\end{tabular}
\end{center}
\end{remark}

\begin{remark}[What remains conjectural]
% label removed: rem:what-remains-conjectural
Three aspects of the programme remain conjectural:
\begin{enumerate}
\item The boundary correspondence for non-abelian theories at the
  chain level (beyond the cohomology-level identification of boundary
  states with bar cocycles).  The Costello--Gwilliam framework provides
  the formal structure, but the explicit quasi-isomorphism
  $\Phi \colon \barBch(\cA_b) \xrightarrow{\sim} \Obs^\bullet_\partial$
  has been proved only in the free case
  (Theorem~\ref{thm:abelian-boundary-correspondence}).
\item The DS reduction functor on bar-cobar data for arbitrary nilpotent
  type (beyond hook-type in type A).  The obstruction is the
  non-complete-intersection property of the BRST complex for
  non-hook-type nilpotents.
\item The BV partition function at genus $g \geq 2$ for the Virasoro
  algebra, where the infinite shadow obstruction tower contributes non-trivially.
  The formula \eqref{eq:virasoro-partition-function} is explicit but the
  convergence of the shadow-tower contributions to the integral over
  $\ov{\cM}_g$ requires the analytic completion programme.
\end{enumerate}
These three open directions are the natural targets for the next stage
of the programme.
\end{remark}
